\section{The complete fixed-shift interface}
\label{sec:literal-interface}

We restate the part of the TPC--61/TPC--68 interface needed below.  The
purpose is to ensure that every later matrix entry is attached to an actual
physical atom and a complete output key.

\subsection{Atoms and recovered arithmetic labels}

Fix a scale \(X\), a nonzero integer \(h_0\), all direct sides and external
blocks in the declared scope, all source and opposite-row blocks, residue
and orbit classes, masks, profiles, endpoint restrictions, and the
very-high cutoff before choosing a coefficient vector.  Let
\(\cW_M\) be the resulting complete set of nonzero missing-high atoms and
let
\[
 \cK_X=\ell^2(\cM_X)
\]
be the full declared source-row space.

Every \(w\in\cW_M\) canonically recovers
\begin{equation}
 m(w)=d_{m(w)}\ell_{m(w)},\qquad
 j(w),\qquad
 n(w)=m(w)j(w)+h_0,
 \label{eq:recovered-row-time-target}
\end{equation}
its largest prime \(p(w)>y\), its cofactor
\[
 r(w)=\frac{n(w)}{p(w)},
\]
and its residual label
\begin{equation}
 s(w)=d_{m(w)}r(w).
 \label{eq:residual-label}
\end{equation}
The physical support is squarefree and gives
\begin{equation}
 \mu(s(w))=\mu(d_{m(w)})\mu(r(w)).
 \label{eq:residual-mobius-factor}
\end{equation}

The complete native output key is denoted by \(i(w)\).  It includes the
direct side, external block, opposite row, \(\gamma\)-coordinate, physical
orbit time \(j\), and every remaining label that is orthogonal in the
native range.  In particular,
\begin{equation}
 i(w)=i(w')\quad\Longrightarrow\quad j(w)=j(w'),
 \label{eq:key-fixes-time}
\end{equation}
but equality of \(i\) does not force equality of \(s\).  The critical-cutoff
singleton property states
\begin{equation}
 \#\{w\in\cW_M:(i(w),m(w))=(i,m)\}\le1.
 \label{eq:im-singleton}
\end{equation}
It does not say that one native key supports only one row, one prime, or
one residual label.

The literal atom driven by a row coefficient \(\zeta_m\) is
\[
 \beta_w\zeta_{m(w)}.
\]
The fixed nonzero coefficient \(\beta_w\) retains every physical source
shape, mask, multiplier, opposite-row factor, output phase, profile, and
terminal weight.  No atomwise phase is introduced after the carrier is
formed.

\subsection{Atomic, residual, and native maps}

Absorb the physical high character in the TPC--61 convention by putting
\begin{equation}
 c_w=-\mu(s(w))\beta_w,\qquad
 (L_X\zeta)_w=c_w\zeta_{m(w)}.
 \label{eq:atomic-lift}
\end{equation}
Define
\begin{align}
 (T_X\zeta)_{i,s}
 &:=
 \sum_{\substack{w\in\cW_M\\i(w)=i,\ s(w)=s}}
 c_w\zeta_{m(w)},
 \label{eq:residual-map}\\
 (\mathsf Ez)_i&:=\sum_s z_{i,s}.
 \label{eq:erasure-map}
\end{align}
Then the literal missing-native map is
\begin{equation}
 \boxed{M_X=\mathsf ET_X.}
 \label{eq:missing-native-factorization}
\end{equation}
Only \(s\) is erased.  Applying a second Möbius sign after
\eqref{eq:erasure-map} would change the physical map.

For every declared row put
\begin{equation}
 W_m:=W_M(m)
 :=
 \sum_{\substack{w\in\cW_M\\m(w)=m}}|\beta_w|^2,
 \qquad
 D_M:=\diag(W_m)_{m\in\cM_X}.
 \label{eq:raw-row-mass}
\end{equation}
This is raw atomic mass, not an average over output fibers.  Moreover
\[
 D_M=L_X^*L_X.
\]
The active missing-row space is
\begin{equation}
 \cM_M:=\{m\in\cM_X:W_M(m)>0\},
 \qquad
 S_M:=\Ran D_M=(\Ker D_M)^\perp
     =\ell^2(\cM_M).
 \label{eq:active-row-space}
\end{equation}
Rows in \(\Ker D_M\) remain present in the full coefficient space and in
the represented map.  Only the normalized missing Gram is restricted to
\(S_M\).

Assume \(S_M\ne\{0\}\), and define
\begin{equation}
 F:=M_XD_M^{\dagger/2}\big|_{S_M},
 \qquad
 F_{\rm res}:=T_XD_M^{\dagger/2}\big|_{S_M}.
 \label{eq:normalized-incidences}
\end{equation}
TPC--61 gives Hermitian zero-diagonal matrices
\[
 R:=H_{\rm res},\qquad E:=H_{\rm erase}
\]
such that
\begin{align}
 F_{\rm res}^*F_{\rm res}&=I+R,
 \label{eq:residual-gram}\\
 F^*F&=I+R+E.
 \label{eq:native-gram}
\end{align}
Thus \(R\) contains common-\((i,s)\) collisions already visible before
erasure, whereas \(E\) contains exactly the unequal-\(s\) cross terms
created at a common \(i\) by \(\mathsf E\).

\subsection{The source-row gauge}

On the squarefree active support, let
\begin{equation}
 U_d=\diag(\mu(d_m))_{m\in\cM_M}.
 \label{eq:row-gauge}
\end{equation}
It is a self-adjoint unitary.  Conjugating the normalized Gram by \(U_d\)
does not change its spectrum, norms, traces of powers, or collision graph.
The signed amplitude of an atom in the gauged normalized incidence is
\begin{equation}
 f_i(m)
 :=
 \frac{\mu(n(w))\beta_w}{\sqrt{W_m}}
 \quad\text{if }w\text{ is the unique atom with }
 (i(w),m(w))=(i,m),
 \label{eq:gauged-key-amplitude}
\end{equation}
and \(f_i(m)=0\) when no such atom exists.  Indeed,
\[
 \mu(d_m)c_w
 =-\mu(d_m)\mu(s(w))\beta_w
 =-\mu(r(w))\beta_w
 =\mu(n(w))\beta_w.
\]

Consequently, for \(m\ne n\), the gauged full defect has entry
\begin{equation}
 (U_d(R+E)U_d)_{mn}
 =
 \sum_i\overline{f_i(m)}f_i(n).
 \label{eq:gauged-full-entry}
\end{equation}
Every summand uses one complete native key.  In the remainder we work in
these gauged coordinates and suppress \(U_d\) from the notation.

\subsection{Direct-cell specialization}

Only in the identical-support direct-cell specialization, a displayed
left-side atom may be written \(w=(m,j,\gamma)\) with
\begin{equation}
 \beta_{m,j,\gamma}
 =a_m(j)K_{m\gamma}q_\gamma(j),
 \qquad
 K_{m\gamma}\in\{0,1\}.
 \label{eq:direct-cell-factorization}
\end{equation}
For two rows define the literal orthogonal output sum
\begin{equation}
 \Gamma_{mn}(j)
 =
 \sum_\gamma
 K_{m\gamma}K_{n\gamma}|q_\gamma(j)|^2.
 \label{eq:Gamma-overlap}
\end{equation}
The \(\gamma\)-coordinate has not been deleted in this formula; it has been
summed exactly in the Gram over orthogonal output coordinates.  Outside the
declared specialization, one must keep the full atom pair sum or place the
complete \(\gamma\)-dependent first-failure indicator inside
\eqref{eq:Gamma-overlap}.

The opposite direct side is orthogonal at the output level.  Hence its Gram
adds to the displayed-side Gram.  This eliminates \(F_L^*F_R\), but it does
not imply
\[
 H^LH^R=0.
\]
An operator walk may therefore use a different side, block, \(\gamma\), or
time on each edge.

\subsection{Six-item physical audit}

The interface satisfies:
\begin{enumerate}[label=\textup{(L1.\arabic*)},leftmargin=4em]
 \item every atom uses the literal coefficient
       \(-\mu(s(w))\beta_w\);
 \item every target is \(m(w)j(w)+h_0\) for one fixed \(h_0\ne0\);
 \item \(D_M\) is raw atomic mass, with no fiber normalization;
 \item source rows and complete native outputs are canonical, and only
       \(s\) is erased;
 \item the carrier and all key partitions are fixed before choosing a
       coefficient vector or an operator extremizer; and
 \item the identities in this paper have zero structural exponent cost,
       while every later estimate remains subject to the strict \(1/400\)
       endpoint ledger.
\end{enumerate}
